For a prime $\mathfrak{q}\subseteq B$, let $\mathfrak{p}$ be its contraction to $A$. It suffices by \exref{5.23} to show $J(B/\mathfrak{q})=0$. Replace $A,B$ by the domains $A/\mathfrak{p}\subseteq B/\mathfrak{q}$; the new $A$ is Jacobson, so $J(A)=0$.

(i) Suppose $B$ is integral over $A$. By \exref{5.5}, $J(B)\cap A=0$. If $0\neq b\in J(B)$, choose a monic equation for $b$ of least degree over $A$. Its constant coefficient is nonzero, since otherwise cancellation of $b$ in the domain $B$ gives a monic equation of smaller degree. But the equation writes that coefficient as $-b$ times an element of $B$, so it lies in $J(B)\cap A=0$, a contradiction.

(ii) If $B$ is finitely generated over $A$, apply \exref{5.22}.

Fields are Jacobson. The ring $\mathbb{Z}$ is Jacobson because its nonzero primes are maximal and the intersection of these maximal ideals is zero. A finitely generated ring is a finitely generated $\mathbb{Z}$-algebra, so the stated consequences follow.
